\documentclass[10pt,a4paper]{amsart}
\usepackage[margin=19mm]{geometry}
\usepackage{amsmath,amssymb,mathtools}
\usepackage[scr=rsfs,cal=euler]{mathalfa}
\usepackage{xfrac}
\usepackage[unicode,hidelinks,bookmarks=false]{hyperref}
\newcommand{\ie}{i.e.\ }
\newcommand{\cf}{cf.\ }
\newcommand{\resp}{respectively\ }
\newcommand{\Subsection}{\subsection}
\providecommand{\nobreakdash}{\nobreak\hbox{-}}
\setlength{\emergencystretch}{2em}
\multlinegap0pt
\usepackage{xcolor}
\usepackage{amssymb}
\usepackage{tikz}
\usepackage{enumerate}
\usepackage[shortlabels]{enumitem}
\usepackage{mathtools}
\newtheorem{lem}{Lemma}
\newtheorem{defi}{Definition}
\newtheorem{thm}{Theorem}
\def\CC{{\mathbf{C}}}
\def\Ga{\Gamma}
\DeclareMathOperator{\Opp}{Opp}
\DeclareMathOperator{\R}{\mathbb{R}}
\DeclareMathOperator{\Th}{Th}
\def\al{\alpha}
\def\ba{\mathbf{a}}
\def\bu{\mathbf{u}}
\def\bv{\mathbf{v}}
\def\bx{\mathbf{x}}
\def\cF{{\mathscr{F}}}
\def\cL{{\mathscr{L}}}
\def\cO{{\mathscr{O}}}
\def\cW{{\mathscr{W}}}
\def\eps{\epsilon}
\def\ga{\gamma}
\def\ol{\overline}
\title{Equidistribution of currents\\ under Anosov group actions}
\author{Nguyen-Bac Dang, Michael Kapovich, Mikhail Lyubich, Shengyuan Zhao}
\date{Source: arXiv:2607.22920v1 (CC BY 4.0)}
\begin{document}
\maketitle

\noindent\emph{Source and licence.} Equidistribution of currents\\ under Anosov group actions. Version arXiv:2607.22920v1 (CC BY 4.0), distributed by arXiv under the Creative Commons Attribution 4.0 International licence, http://creativecommons.org/licenses/by/4.0/, as recorded in the arXiv licence metadata for this exact version and shown on the abstract page https://arxiv.org/abs/2607.22920v1. Full licence notice: \texttt{licenses/arxiv-2607.22920-CC-BY-4.0.txt}.\par\smallskip

\noindent\textit{Editorial scope.} Counted unit: the article's thm contraction (source0.tex lines 12459--12476). The article's complete original proof of it is delivered with it (source0.tex lines 12477--12682, one \texttt{proof} environment of 206 lines). No mathematics is written, repaired or re-derived: the statement and the proof are the article's own lines, in the article's own notation, and no retained line is substituted or re-flowed.  Retained uncounted as the source-local context the proof needs: lines 6773--6773; lines 12170--12172; lines 12250--12254; lines 12261--12261; lines 12269--12279; lines 12301--12317; lines 12403--12412.  Nothing was omitted from any retained span, so the package carries no omission record.  No retained line contains a citation, so the wrapper carries no bibliography and no bibliography entry.  No retained line contains a figure environment, an image-inclusion command or drawing code, so no image is delivered and the package declares no asset dependency.  Editorial formatting only, in addition: an AMS \texttt{amsart} preamble that restates the article's own declarations and macros used by the retained lines, namely the environments \texttt{lem}, \texttt{defi}, \texttt{thm} on separate counters, together with the standard packages those lines need.  The retained environments print in this document as Lemma 1, Definition 1, Theorem 1 in order of appearance, whereas the article prints the same items with the numbering of its own full document.  The complete native source of the article as distributed in the arXiv e-print for this exact version is bundled unchanged as \texttt{sources/205985/source0.tex}.
\subsection{Thickenings in the flag manifold}\label{sec:Thickenings in the flag manifold}

\begin{equation}\label{eq:pm-proj} 
\pi_{w}^{\pm}: \hat{\cO}_{w, \hat{m}}\to \hat{\cL}^\pm_{w,\hat{m}},
\end{equation}

\begin{lem}\label{lem:disjoint}
$\cL^+_{w,m}$ is disjoint from $\bigcup_{v>w} \cO_{v,m}$. 
%2. $\cL^-_{w,m}\setminus s_w(m)\subset \bigcup_{v>w} \cO_{v,m}$. 
%[Maybe even $\cO_{w,m}\setminus \cW^+_w(m)\subset \bigcup_{v>w} \cO_{v,m}$.]  
\end{lem}

\section{Tubes around thickenings}\label{sec:tubes}

\begin{defi}\label{def:tube}
Given such $\CC$,  
we define the {\em $\CC$-tube} around $S_{w^\vee}(x_-)$, denoted $N_{\CC}(S_{w^\vee}(x_-))$, as
$$
\exp(\CC)U^-_w \cdot x_w. 
$$
Similarly, when $\bx = \beta^2(\hat{m})$ with $\hat{m} \in \hat\Gamma$, we set $N_C(\hat{m} \times S_{w^\vee}(x_-))$ to be the set: 
\begin{equation*}
\exp(\CC) U_w^- \cdot (\hat{m},x_w )= \{ \hat{m} \} \times N_{\CC}(S_{w^\vee}(x_-))  \subset \hat{E}_{\hat{m}}. 
\end{equation*}
\end{defi} 

\begin{lem}\label{lem:compactifications of leaves} 
1. The partial boundary $\partial' N_{\CC}(S_{w^\vee}(x_-))$ of $N_{\CC}(S_{w^\vee}(x_-))$ in $\cF$ equals
$$
\partial' S_{w^\vee}(\exp({\CC})x_-)= \bigcup_{g\in \exp({\CC})} g \partial' S_{w^\vee}(x_-)= \bigcup_{g\in \exp({\CC})}  \partial' S_{w^\vee}(g x_-). 
$$
 In particular, $\partial' N_{\CC}(S_{w^\vee}(x_-))$ is compact and equals the thickening $\Th_{<w^\vee}( \exp({\CC})x_-)$, where 
$$
 \Th_{< u}= \{v\in W: v< u\}. 
$$ 

2. When $\bx = \beta(\hat{m})$ and $m=\Pi(\hat{m})$,
the projection of $\{ \hat{m}\} \times \partial'(N_{\CC}(S_{w^\vee}(x_-)))$ to $E_m$ is contained in $\bigcup_{v>w} \cO_{v,m}$ and, thus,
% $\partial' N_{\CC}(S_{w^\vee}(x_-))\subset \bigcup_{v>w} \cO_{x_v}$ and, thus,  
$$
\ol{N_{\CC}(\{ \hat{m} \} \times S_{w^\vee}(x_-))}= N_{\CC}(\{ \hat{m} \}\times \Th_{w^\vee}(x_-))\subset \bigcup_{v\ge w} \hat{\cO}_{v, \hat{m}}. 
$$
\end{lem}

\begin{lem}\label{lem:comparison}
$||\cdot||^w$ is a measurable Finsler metric on $T\hat{\cO}_{\hat{m}}^w$ which 
is locally bilipschitz equivalent to the restriction of the Riemannian metric on $\cF$ in the following sense:  

There exists a positive continuous function $c_w: \hat{\cO}_{\hat{m}}^w\to \R_+$  such that:
 $$
\frac{1}{c_w(z)}||\bv|| \leqslant ||\bv||^w\leqslant c_w(z)||\bv||, 
$$
where  $z\in \hat{\cO}_{\hat{m}}^w$ and where $\bv\in T_z\hat{\cO}_{\hat{m}}^w$ is a vertical vector.  
\end{lem}

\begin{thm}\label{thm:transvection contraction} 
%\textcolor{red}{
%Fix $\eps>0$. Then for every $w\in W$ there exists a pair of continuous functions $C_w^\pm: \cO_{x_w}\to \R_+$ 
%such that the following holds for every unit vector $\ba\in \mathfrak{a}^+_{\eps}$, every $z\in \cO_{x_w}$ and $t\geqslant 1$: 
%$$
%||de^{\mp t\ba}||_z^\pm \leqslant C_w^\mp(z) e^{-\eps t}.  
%$$
%Furthermore, $C_w^+$ is bounded on each tube $N_{\CC}(\Th_w(x_+))$ and $C_w^-$ is bounded on each tube $N_{\CC}(\Th_{w^\vee}(x_-))$. In particular, 
%$C_w^\pm$ is bounded along each semi-orbit of $e^{\mp t\ba}$, $t\geqslant 0$. }
Fix $\bx = \beta^2(\hat{m})$ with $\hat{m} \in \partial_\infty^2 \Ga \setminus \Delta$. 
and  $\eps>0$. Then for every $w\in W$ there exists a pair of continuous functions $C_w^\pm: \hat{\cO}_{w,\hat{m}}\to \R_+$ 
such that the following holds for every unit vector $\ba\in \mathfrak{a}^+_{\eps}$, every $z\in \hat{\cO}_{w,\hat{m}}$ and $t\geqslant 1$: 
$$
||de^{\mp t\ba}||_z^\pm \leqslant C_w^\pm(z) e^{-\eps t}.  
$$
Furthermore, $C_w^+$ is bounded on each tube $N_{\CC}(\Th_w(x_-))$ and $C_w^-$ is bounded on each tube $N_{\CC}(\Th_{w^\vee}(x_+))$. In particular, 
$C_w^\pm$ is bounded along each semi-orbit of $e^{\mp t\ba}$, $t\geqslant 0$. 
\end{thm}

\begin{proof} We will prove the claim for the action $e^{-t\ba}$ on the direction induced by the foliation $\hat{\cW}^+_w$ as the claim for  $e^{t\ba}$ follows by reversing the roles of $x_+, x_-$ and exchanging the roles of $\hat{\cW}^+_w$ with $\hat{\cW}^-_w$. 

Note that the contraction estimate is vacuous in the case $w=1$ since then there are no nonzero $w$-horizontal tangent vectors. 
(Tangent vectors to leaves of the horizontal foliation $\cW^+_{x_w}$ on $\cO_{x_w}$ are $w$-horizontal.) 


We equip ${\mathfrak g}$ with a hermitian inner product with respect to which the root subalgebras ${\mathfrak g}_\al, \al\in \Phi$, 
are pairwise orthogonal. We let $||\cdot||_0$ denote the corresponding norm. 

\medskip 

We first prove the contraction estimate for the Finsler metric:
\begin{equation}\label{eq:contract} 
 ||d\ga (\bv)||^w\leqslant  e^{-\eps t} ||\bv||^w 
\end{equation}
for every $w\in W$ and every $w$-horizontal vector $\bv$ in $\hat{\cO}_{w,\hat{m}}$, where $\gamma = e^{-t \ba}$. 
The proof is by induction on $|w^\vee|=n-|w|$. 

\medskip 
{\bf Base of induction.} Suppose that $w=w_0, w^\vee=1$. Then $\hat{\cO}_{w,\hat{m}}=\Opp(x_+)$ and the horizontal foliation $ \hat{\cW}_{w,\hat{m}}^+$ consists of a single leaf, namely, $\hat{\cO}_{w,\hat{m}}$. 

 
 Recall that the condition $\ba\in \mathfrak{a}^+_{\eps}$ means that for every positive root $\al\in \Phi^+$ we have
 $\al(\ba)\geqslant \eps$. But the numbers  $e^{\al(t\ba)}$ are the eigenvalues of $Ad(e^{t\ba})$ acting on ${\mathfrak u}$, $t\in \R$, where ${\mathfrak u}$ is the Lie algebra of the unipotent radical of the stabilizer of $x_+$.  Hence, for the norm $||\cdot||_w$, the contraction estimate at $x_w=x_-$ for a vector $\bv \in T_{(\hat{m}, x_-)} \hat{\cO}_{w, \hat{m}}$ becomes a computation on ${\mathfrak u}$, which is 
 $$
||de^{-t\ba}(\bv)||_0= ||de^{-t\ba}(\bv)||^w \leqslant e^{-\eps t} ||\bv||_0= e^{-\eps t} ||\bv||^w. 
 $$
 
For general points $z\in \hat{\cO}_{w, \hat{m}}$ we get a contraction estimate as 
follows. Let $\bv\in T_z \hat{\cO}_{w, \hat{m}}$, $\ga=e^{-t\ba}$, $\ba$ is a unit vector in $\mathfrak{a}_\eps^+$. Our first goal is to compare the norms 
$$
||\bv||_w, \quad ||d\ga (\bv)||_w. 
$$
Let $g, h\in U_{x_w}$ be such that $g(x_-)=z$, $h(\ga z)=x_-$. 

\begin{lem}
 $h\circ \ga \circ g=\ga$. 
\end{lem} 
\begin{proof} By the choice of $g, h$, the composition $\hat\ga:=h\circ \ga \circ g$ fixes the point $x_-$. Note that $\hat\ga$ belongs to 
$B_{w^\vee}= U_w\rtimes T$ (see \S~\ref{sec:Thickenings in the flag manifold}). Since $T$ fixes $x_w=x_-$ and $U_w$ acts simply transitively on $\cO_{x_w}$, the stabilizer of $x_w=x_-$ in 
$B_{w^\vee}$ equals $T$. Hence, $\hat\ga\in T$. Since $U_w$ is a normal subgroup of $B_{w^\vee}$, we also have 
$$
\hat\ga:=h\circ \ga \circ g= \ga \circ h^\ga \circ g, h^\ga \circ g\in U_w. 
$$
%Under the quotient map $B_{w^\vee}\to B_{w^\vee}/U_w\cong T$, the images of $\ga$ and $\hat\ga$ are 
%the same. Hence, 
Since $\hat\ga\in T$, it follows that $\hat\ga=\ga$, as claimed. 
\end{proof}


Since the metric 
$||\cdot||_w$ is $U_{x_w}$-invariant, it follows that for $\bu= dg^{-1}(\bv)$, 
$$
||\bu||_0=||\bu ||_w= ||\bv||_w, \quad  ||d\ga (\bv)||_w= ||d\ga (\bu)||_0. 
$$

Hence, we obtain the same exponential contraction estimate 
$$
||de^{-t\ba}(\bv)||^w \leqslant e^{-\epsilon t}  ||\bv||^w 
 $$
 at all points $z\in \hat{\cO}_{w, \hat{m}}$ 	and all $\bv \in T_z \hat{\cO}_{w, \hat{m}}$. This concludes the proof of \eqref{eq:contract} in the case $w=w_0$. 





\medskip 
{\bf Inductive step.} We assume that \eqref{eq:contract} holds for all $w\in W$ with $n-|w|<k$. Consider $w$ with $n-|w|=k$.  

%Using $U_w$-invariance of the metric $||\cdot||_w$, we first check the horizontal contraction at $x_w$, then (using induction) 
%at points $\cL_w^-$ (with respect to $||\cdot||^w$), then at general leaves of $\cW^-_w$  
%using Lemma \ref{lem:compactifications of leaves} %(via splitting $e^{t\ba}$ as a composition). 
%Then conclude using Lemma \ref{lem:comparison}. 

Arguing as in the base of induction case, and taking into account the fact that action of $U_{x_w}$ 
preserves the vertical/horizontal foliations of $\hat{\cO}_{w,\hat{m}}$,   the contraction estimate 
\begin{equation}\label{eq:w-estimate} 
||d\ga(\bv)||_w\leqslant e^{-\eps t} ||\bv||_w,
\end{equation}
(for $w$-horizontal vectors $\bv\in T \hat{\cO}_{w,\hat{m}}$)  reduces to the case of the base-point $z= (\hat{m},x_w)$. But then the claim follows from the fact that the eigenvalues 
of $Ad(\ga)$ acting on ${\mathfrak u}^+_w$ are given by $e^{-t\al(\ba)}$ for $\al\in \Phi_w^+\subset \Phi^+$. By the assumption, all these eigenvalues are 
$\leqslant  e^{-\eps t}$, as required. 

In order to get an estimate in terms of $||\cdot||^w$, we have to consider two cases:  

{\bf Case 1.} Assume that $z\in \hat{\cO}_{w,\hat{m}}$ and $\bv\in T_z \hat{\cO}_{w,\hat{m}}$ are such that  
$$
||\bv||_v= ||\bv||^w, \quad ||d\ga(\bv)||_v= ||d\ga(\bv)||^w 
$$
for some $v\in W, v \geqslant w$. If $v = w$, then the inequality \eqref{eq:contract} follows from \eqref{eq:w-estimate}. Otherwise, $v> w$ and as $n- |v| < n-|w|$, the induction assumption can be applied to $v$ and we conclude.
% as $n- |v| < n- |w|=k$. 
For instance, this applies 
to all points $z\in \{ \hat{m} \} \times S_{w^\vee}(x_-)=\cL_{w,\hat{m}}^-$, since $\cL_{w, \hat{m}}^-$ is disjoint from 
$$
\bigcup_{v <  w} \hat{\cO}_{v, \hat{m}}
$$
(see Lemma \ref{lem:disjoint} with the sign reversed, exchanging $x_+$ with $x_-$ and $\mathcal{L}^+$ with $\cL^-$) and is $T$-invariant. Hence,  for all vectors $\bv$ tangent to $\hat{\cW}_w^+$ at points of $\cL_w^-$ we have
$$
||\bv||_w=||\bv||^w. 
$$
%We now consider horizontal tangent vectors $\bv$ at points $z\in \cO_{x_w}\setminus \cL_{x_w}^-$. 

\medskip 
{\bf Case 2.} 
$$
||\bv||^v= ||\bv||_v= ||\bv||^w, \quad ||d\ga(\bv)||^w= ||d\ga(\bv)||_{v'} 
$$
for some $ v \neq  v'\in W$ satisfying $v\ge w, v'\ge w$.
Observe that this implies that $z \in \hat{\cO}_{wv,\hat{m}}\cap \hat{\cO}_{wv',\hat{m}}$ since $T$ preserve each domain $\hat{\cO}_{u}$ for $u \in \{ w,v,v'\}$. 
So that the vector $\bv$ belongs to $\hat{\cW}_{v,\hat{m}}^+ \cap \hat{\cW}_{w,\hat{m}}^+ \cap \hat{\cW}_{v', \hat{m}}^+$. 
If $v' > w$ then by the induction hypothesis and using the fact that $|| \cdot ||^{v'}\leqslant || \cdot ||^{w}$, we have: 
\begin{equation*}
|| d \gamma(\bv)||^w = || d\gamma (\bv)||^{v'} \leqslant e^{-\epsilon t} || v||^{v'} \leqslant e^{-\epsilon t} || v||^{w},  
\end{equation*}
as required. 
Otherwise, we have $v'=w$ and using \eqref{eq:w-estimate} 
we obtain: 
\begin{equation*}
|| d \gamma(\bv)||^{w} = || d\gamma(\bv) ||_w \leqslant e^{-\epsilon t} || \bv ||_w \leqslant e^{-\epsilon t} || \bv||^w, 
\end{equation*}
as required.   
%Observe that \eqref{eq:w-estimate} shows that the norm $|| \cdot ||_u$ for each $u \in W$ decays exponentially. 
% Then there exists $T\in [0,t]$ such that for 
%$$
%\bv':= de^{-\ba T}(\bv),  ||\bv'||^v=  ||\bv'||^{v'}. 
%$$
%By applying Case 1 to the intervals $[0,T]$ and $[T,t]$ (for the second case, we apply Case 1 to the vector $\bv'$) and taking into account the fact that 
%$\bv$ is $v$-horizontal and $\bv'$ is $v'$-horizontal, we obtain:
%$$
%||\bv'||_v\leqslant e^{-\eps T}||\bv||_v, $$
%\begin{align*}
% ||d\ga(\bv)||_{v'}= ||de^{-\ba (t-T)}(\bv')||_{v'}\leqslant % ||d\ga(\bv)||_{v'}\leqslant 
%\\
%e^{-\eps (t-T)} ||\bv'||_{v'}=e^{-\eps (t-T)} ||\bv'||_v.
%\end{align*}
%Combining these inequalities, we get 
%$$
% ||d\ga(\bv)||^w\leqslant e^{-\eps t}  ||d\ga(\bv)||^w,
%$$
%as required. 
This proves the inequality 
$$
 ||d\ga (\bv)||^w\leqslant  e^{-\eps t} ||\bv||^w  
$$
at all points $z\in \hat{\cO}_{w,\hat{m}}$. 

\medskip 
{\bf Estimating contraction in terms of the Riemannian metric on $\cF$.}  
The restriction of 
the Riemannian norm on $T\cF$ to any compact in $\cO^w_{\bx}$  is uniformly bilipschitz on compacts to the measurable Finsler metric  
$||\cdot||^w$ (cf. Lemma \ref{lem:comparison}). The issue is that even if $z\in \cO^w_{\bx}$
 lies in a given compact subset, we, a priori, do not know where its image under $\ga$ is. Consider 
 $z\in \cO_{x_w}$. Recall that earlier in \eqref{eq:pm-proj} we defined continuous maps $\pi_{w}^{\pm}: \cO_{x_w}\to \cL^\pm_{w,\bx}$. 
 Set $z_\pm:= \pi_{w}^{\pm}(z)$  and define 
 $$
 u_z:= \Psi_w(z_+)\in \mathfrak{u}_w^+\subset 
 \mathfrak{u}_w, \quad R=R(z):= ||u_z||_0. 
 $$
 %Define $z_+\in S_w(x_+)$ as the unique intersection point of $U_w^- z$ and $S_w(x_+)$. 
 %This point depends continuously on $z$. 
 % Let $u_z:= \Psi_w(z_+)\in \mathfrak{u}_w^+\subset \mathfrak{u}_w$. Set $R:= ||u_z||_0$. 
 Thus, $R(z)$ is a continuous function of $z$ and $z$ belongs to the $\CC$-tube $N_{\CC}(S_{w}(x_-))\subset \cO_{x_w}$, where 
$$
\CC=\bar{\mathbf{B}}(0, R)\subset \mathfrak{u}^+_w
$$ 
is the closed $R$-ball with respect to the norm $||\cdot||_0$ on $ \mathfrak{u}_w$. 
(See Definition \ref{def:tube} for the notion of $\CC$-tubes.) 
Clearly, $\CC$ is a compact convex subset of $\mathfrak{u}^+_w$ containing $0$. 

We proved in Lemma \ref{lem:compactifications of leaves} 
that the closure in $\cF$ of $N_{\CC}(S_{w}(x_-))$ equals the compact 
%$$
%N_{\CC}(\Th_{w}(x_-))\subset \cO_{\bx}^w$$

{$$
N_{\CC}(\{\hat{m} \} \times \Th_{w}(x_-))\subset \bigcup_{v\geq w^\vee}\hat{\cO}_{v, \hat{m}}=\hat{\cO}_{\hat{m}}^{w^\vee}.$$}
 
Moreover, in \S \ref{sec:tubes} we also observed that 
for every $\gamma\in \exp(- \mathfrak{a}^+)$, 
$$
\ga(N_{\CC}(\Th_{w}(x_-)))\subset N_{\CC}(\Th_{w}(x_-)).$$
Define the function $C_w^-(z)$ as follows.

%\textcolor{red}{
%Recall that in Lemma \ref{lem:comparison} we proved existence of a positive continuous function $c_{w^\vee}$ 
%on $\cO^{w^\vee}_{\bx}$ providing bilipschitz comparison between the norms $||\cdot||$ and $||\cdot||^{w^\vee}$ on the tangent bundles $T\cO_{\bx}^{w^\vee}$. 
%Then we set
%$$
%C_w^-(z):=\max_{x\in N_{\CC}(\Th_{w}(x_-))} c_{w^\vee}(x). 
%$$
%Continuity of the function $R(z)$ implies continuity of $C_w^-(z)$.
%}

Recall that in Lemma \ref{lem:comparison} we proved existence of a positive continuous function $c_{w^\vee}$ 
on $\hat{\cO}^{w^\vee}_{\hat{m}}$ providing bilipschitz comparison between the norms $||\cdot||$ and $||\cdot||^{w^\vee}$ on the tangent bundles $T\hat{\cO}_{\hat{m}}^{w^\vee}$. 
Then we set
$$
C_w^+(z):=\max_{x\in N_{\CC}(\Th_{w}(x_-))} c_{w^\vee}(x). 
$$
Continuity of the function $R(z)$ implies continuity of $C_w^+(z)$.
 
We, thus, obtain (using the inequality \eqref{eq:contract}) that for any $\bv \in T_z \hat{\cO}_{w,\hat{m}} \cap \hat{\cW}^+_{\hat{m}}$, one has:
\begin{equation*}
|| d (e^{-t \ba}) \bv ||_{e^{-t  \ba } \cdot z} \leqslant C_w^+(z) || d (e^{-t \ba}) \bv ||^w \leqslant C_w^+(z) e^{-\epsilon t} ||\bv ||^w \leqslant (C_{w}^+(z))^2 e^{-\epsilon t} ||\bv ||_{z}, 
\end{equation*}
Theorem follows. \end{proof}

\end{document}
